% ============================================================================
% Chapter 3 — Mathematical Formulation
%
% Equation-numbering contract with the companion code (do not renumber):
%   (3.2) impact, (3.5) vitality ODE, (3.6) capital ODE, (3.7) complexity
%   coupling, (3.14) carrying capacity, (3.15) comparative statics.
% Definition 5 = depletion ratio; Definition 7 = strategic threshold.
% All other displayed mathematics in this section must remain unnumbered.
% ============================================================================
\section{Mathematical Formulation}
\label{sec:formulation}

\subsection{State variables}
\label{sec:state}

\begin{definition}[Subjective vitality]
\label{def:vitality}
$V(t)$ denotes the leader's subjective vitality at time $t$: the energetic
resource available for deliberate leadership work
\citep{ryan1997subjective,meijman1998effort}. Vitality is measured in
arbitrary energy units and confined to $[0,\Vmax]$, where $\Vmax$ is the
individual's recovery ceiling.
\end{definition}

\begin{definition}[Career capital]
\label{def:capital}
$C(t)\ge 0$ denotes accumulated career capital: the stock of human, social,
and reputational assets that confer status and opportunity
\citep{becker1993human}.
\end{definition}

The state space is
\begin{equation}
\label{eq:statespace}
(V(t),\,C(t)) \in [0,\Vmax]\times[0,\infty),
\end{equation}
and Appendix~\ref{app:invariance} verifies that the flow defined below never
leaves it (positive invariance). Time is measured in weeks; one unit is of
the order of one working week, though the calibration is illustrative
(\S\ref{sec:calibration}).

\subsection{Functional forms}
\label{sec:forms}

\begin{definition}[Organizational complexity]
\label{def:complexity}
$O(t) \ge O_0 > 0$ denotes the organizational complexity load facing the
leader: the volume of interfaces, coordination demands, and matters requiring
attention \citep{garicano2000hierarchies,hambrick2005executive}. $O_0$ is the
irreducible baseline load; the coupling of $O$ to the leader's own capital is
specified in \eqref{eq:complexity} below.
\end{definition}

\begin{definition}[Energy-gated impact]
\label{def:impact}
Leadership impact is
\begin{equation}
\label{eq:impact}
I \;=\; \frac{C\,V}{1+\varphi O},
\end{equation}
with suppression coefficient $\varphi>0$.
\end{definition}

Impact requires both assets ($C$) and energy ($V$)---the energetic-deployment
constraint absent from classical human-capital theory---and is attenuated by
complexity friction in the denominator: as $O\to\infty$, $I\to 0$ for any
finite $(V,C)$.

Vitality dynamics balance two flows. Recovery follows a saturating
(logistic-type) law toward the ceiling,
\begin{equation}
\label{eq:recovery}
\text{recovery}(V) \;=\; R\Bigl(1-\frac{V}{\Vmax}\Bigr),
\end{equation}
with maximal restoration rate $R$ realized at full depletion
\citep{sonnentag2007recovery}. Drain grows as a power law in complexity and is
regularized near $V=0$ by a smooth barrier,
\begin{equation}
\label{eq:drain}
\text{drain}(V,O) \;=\; \delta\,O^{\gamma}\,\frac{V}{V+\varepsilon},
\qquad \gamma \ge 1,\ \varepsilon>0,
\end{equation}
where $\delta$ is the energetic cost coefficient and $\gamma$ the complexity
exponent. The factor $V/(V+\varepsilon)$ encodes that a fully depleted leader
cannot be drained further (disengagement at exhaustion) and guarantees
positive invariance (Appendix~\ref{app:invariance}); analytical results are
stated for $\varepsilon>0$ with the $\varepsilon\to0$ limit tabulated in
Appendix~\ref{app:epsilon}.

\subsection{The dynamical system}
\label{sec:system}

The DLVT model is the autonomous planar system
\begin{equation}
\label{eq:dVdt}
\frac{\dd V}{\dd t}
 \;=\; R\Bigl(1-\frac{V}{\Vmax}\Bigr)
 \;-\; \delta\,O^{\gamma}\,\frac{V}{V+\varepsilon},
\end{equation}
\begin{equation}
\label{eq:dCdt}
\frac{\dd C}{\dd t} \;=\; \alpha\,I \;-\; \mu\,C,
\end{equation}
where $\alpha$ is the capital accumulation rate and $\mu$ the depreciation
rate, and the loop is closed by the complexity coupling
\begin{equation}
\label{eq:complexity}
O \;=\; O_0 + \beta\,C^{\eta},
\qquad \beta>0,\ \eta\ge 1 .
\end{equation}
Substituting \eqref{eq:impact} into \eqref{eq:dCdt} gives
$\dot C = C\,\bigl(\alpha V/(1+\varphi O)-\mu\bigr)$, so capital dynamics
factor through $C$ itself; this structure carries the scope-absorption result
of \S\ref{sec:equilibrium-lowvitality}.

\begin{definition}[Depletion ratio]
\label{def:depletion}
The depletion ratio is
\begin{equation}
\label{eq:gamma-ratio}
\Gamma(t) \;=\; \frac{\delta\,O(t)^{\gamma}}{R}.
\end{equation}
$\Gamma>1$ indicates that the drain ceiling exceeds the recovery ceiling: no
vitality level near the top of the range can be sustained, and the
quasi-equilibrium of \S\ref{sec:transient} collapses toward the depletion
band $\{V<\varepsilon\}$.
\end{definition}

\begin{assumption}[Scope condition: endogenous load]
\label{ass:scope}
Complexity is a deterministic, memoryless image of the leader's own capital,
\eqref{eq:complexity}. Exogenous load shocks, peer effects, and
organizational restructuring are outside the model's scope; the exogenous
variant $\dot V$ with prescribed $C(t)$ (used in
Proposition~\ref{prop:band}(ii)) is the only departure we analyze.
\end{assumption}

\subsection{Positive invariance}
\label{sec:invariance}

At $V=0$ the drain term vanishes and $\dot V|_{V=0}=R>0$, so $V=0$ is
\emph{repelling}; at $V=\Vmax$, $\dot V<0$; at $C=0$, $\dot C=0$, so the
$C$-axis is invariant and capital cannot become negative. Hence
$[0,\Vmax]\times[0,\infty)$ is forward invariant. The elementary verification
is Appendix~\ref{app:invariance}. Two consequences deserve emphasis. First,
trajectories cannot reach $V=0$: the ``finite-time depletion to zero''
theorem of earlier drafts was false as stated and is replaced by
Proposition~\ref{prop:band}. Second, $\{C=0\}$ being invariant means the
global-stability statement of Theorem~\ref{thm:gas} is necessarily about the
\emph{open} set $\{C>0\}$.

\subsection{Nullcline geometry and transient depletion}
\label{sec:transient}

Setting $\dot V=0$ in \eqref{eq:dVdt} and clearing denominators shows that,
for each fixed complexity level $O$, the vitality nullcline solves the
quadratic
\begin{equation}
\label{eq:vqe-quadratic}
\frac{R}{\Vmax}V^{2}
\;-\;\Bigl(R-\frac{R\varepsilon}{\Vmax}-\delta O^{\gamma}\Bigr)V
\;-\;R\varepsilon \;=\;0 .
\end{equation}

\begin{proposition}[Vitality nullcline: single branch, strictly decreasing]
\label{prop:nullcline}
For every $O>0$ (equivalently every $C\ge0$), the quadratic
\eqref{eq:vqe-quadratic} has exactly one positive root $V_{\mathrm{qe}}(O)$,
and along the nullcline $V_{\mathrm{qe}}$ is strictly decreasing in $C$.
\end{proposition}

\begin{proof}
The product of the roots of \eqref{eq:vqe-quadratic} is
$-\varepsilon \Vmax<0$, so the roots are real of opposite sign: exactly one
positive root. Writing $F(V,C)$ for the right side of \eqref{eq:dVdt},
$\partial F/\partial V = -R/\Vmax - \delta O^{\gamma}\varepsilon/(V+\varepsilon)^{2}<0$
and
$\partial F/\partial C = -\delta\gamma O^{\gamma-1}O'(C)\,V/(V+\varepsilon)<0$
for $V,C>0$; by the implicit function theorem
$\dd V_{\mathrm{qe}}/\dd C = -F_C/F_V<0$.
\end{proof}

The positive root $V_{\mathrm{qe}}(O)$ is the \emph{quasi-equilibrium} of the
vitality equation at frozen complexity. At the baseline calibration with
$O=O_0=1$ it sits near the ceiling, $V_{\mathrm{qe}}\approx 9.93$
($\approx0.99\,\Vmax$); the closed-form approximation
$V_{\mathrm{qe}}\approx R\varepsilon/(\delta O^{\gamma})$ used in earlier
drafts is valid \emph{only} in the strong-drain limit
$\delta O^{\gamma}\gg R$ and is badly wrong at moderate load (at $O=1$ it
returns $15>\Vmax$).

\begin{proposition}[Finite-time band entry; corrected]
\label{prop:band}
Consider the vitality equation \eqref{eq:dVdt} with an exogenous complexity
path $O(t)$ (Assumption~\ref{ass:scope}).
\begin{itemize}
\item[(i)] $V=0$ is repelling: $\dot V|_{V=0}=R>0$. In particular no
trajectory reaches $V=0$ in finite or infinite time, and the ``finite-time
depletion to zero'' claim of earlier drafts is retracted.
\item[(ii)] If $O(t)\ge\Omega$ for all $t\ge t_0$ with
$\Omega^{\gamma} > 2R/\delta$, then any trajectory with
$V(t_0)=V_0\ge\varepsilon$ enters the depletion band
$\{V<\varepsilon\}$ no later than
\begin{equation}
\label{eq:Tstar}
T^{*} \;=\; t_0 \;+\;
\frac{V_0-\varepsilon}{\;\tfrac{1}{2}\delta\Omega^{\gamma}-R\;}.
\end{equation}
\item[(iii)] If instead $O(t)\equiv O$ is constant, $V(t)$ converges
monotonically and exponentially to the positive quasi-equilibrium
$V_{\mathrm{qe}}(O)$, the positive root of \eqref{eq:vqe-quadratic}.
\end{itemize}
\end{proposition}

\begin{proof}[Proof sketch]
(i) is immediate from \eqref{eq:dVdt}. (ii) For $V\ge\varepsilon$ we have
$V/(V+\varepsilon)\ge\tfrac12$, so
$\dot V \le R-\tfrac12\delta\Omega^{\gamma}<0$, a uniformly negative bound;
integrating gives \eqref{eq:Tstar}. (iii) At fixed $O$ the scalar flow has a
single positive rest point by Proposition~\ref{prop:nullcline}, with
$F_V<0$ there, so convergence is exponential at rate $|F_V(V_{\mathrm{qe}})|$.
\end{proof}

Thus sustained overload drives vitality into an $\varepsilon$-neighborhood of
zero in finite time---the phenomenologically relevant statement---while zero
itself is never attained.

\subsection{Equilibrium structure: existence and uniqueness}
\label{sec:equilibrium}

Interior equilibria satisfy $\dot C=0$ with $C>0$, which by
\eqref{eq:impact}--\eqref{eq:dCdt} forces $\alpha V/(1+\varphi O)=\mu$, i.e.\
the \emph{capital nullcline}
\begin{equation}
\label{eq:cnullcline}
V_c(O) \;=\; \frac{\mu}{\alpha}\,\bigl(1+\varphi O\bigr),
\end{equation}
an increasing affine function of $O$. Substituting \eqref{eq:cnullcline} into
the vitality equation defines the reduced residual
\begin{equation}
\label{eq:phi-residual}
\Phi(O) \;=\;
R\Bigl(1-\frac{V_c(O)}{\Vmax}\Bigr)
\;-\;\delta\,O^{\gamma}\,\frac{V_c(O)}{V_c(O)+\varepsilon},
\end{equation}
whose roots with $O>O_0$ (equivalently $C>0$) are exactly the interior
equilibria.

\begin{theorem}[Existence and uniqueness, global]
\label{thm:uniqueness}
For every parameter vector in $\mathbb{R}^{11}_{>0}$, the system
\eqref{eq:dVdt}--\eqref{eq:complexity} has \emph{at most one} interior
equilibrium $(V^{*},C^{*})$ in the entire positive orthant. An interior
equilibrium exists if and only if $\Phi(O_0)>0$, which holds at the baseline
calibration.
\end{theorem}

\begin{proof}
Differentiate \eqref{eq:phi-residual}:
\[
\Phi'(O)
= -\frac{R}{\Vmax}\,V_c'(O)
\;-\;\delta\gamma O^{\gamma-1}\,\frac{V_c(O)}{V_c(O)+\varepsilon}
\;-\;\delta O^{\gamma}\,\frac{\varepsilon\,V_c'(O)}{\bigl(V_c(O)+\varepsilon\bigr)^{2}} .
\]
With $V_c'(O)=\mu\varphi/\alpha>0$, every one of the three terms is strictly
negative for $O>0$: the first because $R,\Vmax>0$; the second because
$V_c>0$; the third because $\varepsilon>0$. Hence $\Phi$ is strictly
decreasing on $(0,\infty)$. Since $O(C)=O_0+\beta C^{\eta}$ is strictly
increasing in $C$, the map $C\mapsto\Phi(O(C))$ is strictly decreasing, so it
has at most one root: at most one interior equilibrium, with no restriction
to any bounded region. For existence: $\Phi$ is continuous, strictly
decreasing, and $\Phi(O)\to-\infty$ as $O\to\infty$ (the drain term diverges
while the recovery term is bounded above), so a root with $O>O_0$ exists iff
$\Phi(O_0)>0$. At baseline, $V_c(O_0)=2.3$ and
$\Phi(O_0)=3(1-0.23)-0.02\cdot(2.3/2.4)\approx 2.29>0$.
\end{proof}

Uniqueness in the \emph{entire} positive orthant is stronger than the
bounded-region statements of earlier drafts and matters for what follows: it
excludes multistability outright, for all positive parameter values, and it
is one of the three ingredients of the global-stability proof
(Appendix~\ref{app:gas}).

\subsection{Baseline calibration}
\label{sec:calibration}

Table~\ref{tab:parameters} reports the baseline calibration used in all
figures and in the companion code (initial capital $C_0=5.0$ throughout).

\begin{table}[t]
\centering
\caption{Baseline calibration (Table~1). \emph{The calibration is
illustrative}: it is chosen for a well-conditioned phase portrait in
dimensionless units and has no empirical anchor. No quantitative claim in
this paper should be read as calibrated to data; qualitative claims are
stated with their parameter dependence. Derived quantities are reported at
$\varepsilon=0.1$; see Appendix~\ref{app:epsilon} for the
$\varepsilon\to0$ extrapolation.}
\label{tab:parameters}
\vspace{0.4em}
\begin{tabular}{clll}
\toprule
Symbol & Code key & Baseline & Interpretation \\
\midrule
$R$            & \texttt{R}     & $3.0$  & vitality recovery rate \\
$\Vmax$        & \texttt{Vmax}  & $10.0$ & maximum vitality capacity \\
$\delta$       & \texttt{delta} & $0.02$ & energetic cost coefficient \\
$\gamma$       & \texttt{gamma} & $2.0$  & complexity exponent in drain \\
$O_0$          & \texttt{O0}    & $1.0$  & baseline organizational complexity \\
$\beta$        & \texttt{beta}  & $0.25$ & capital--complexity coupling \\
$\eta$         & \texttt{eta}   & $1.0$  & capital scaling exponent \\
$\alpha$       & \texttt{alpha} & $0.1$  & capital accumulation rate \\
$\varphi$      & \texttt{phi}   & $0.15$ & complexity--impact suppression \\
$\mu$          & \texttt{mu}    & $0.2$  & capital depreciation rate \\
$\varepsilon$  & \texttt{eps}   & $0.1$  & smooth-barrier regularization \\
\midrule
\multicolumn{4}{l}{\emph{Derived quantities at baseline ($\varepsilon=0.1$)}}\\
$V^{*}$        &                & $4.70$   & equilibrium vitality (\S\ref{sec:equilibrium-lowvitality}) \\
$C^{*}$        &                & $32.03$  & equilibrium capital \\
$O^{*}$        &                & $9.01$   & equilibrium complexity \\
$\lambda_{1,2}$&                & $-0.205\pm0.331i$ & Jacobian eigenvalues (App.~\ref{app:jacobian}) \\
$\Cmax$        &                & $45.0$   & carrying capacity, $\eta=1$ (Prop.~\ref{prop:carrying}) \\
$\Cmax|_{\eta=1.5}$ &           & $12.7$   & carrying capacity at $\eta=1.5$ \\
$\Cmax|_{\eta=2}$   &           & $6.7$    & carrying capacity at $\eta=2$ \\
$\Ctrap$       &                & $102.67$ & trapping bound (Thm.~\ref{thm:gas}); $\Ctrap\neq\Cmax$ \\
\bottomrule
\end{tabular}
\end{table}

Because the calibration is illustrative, we report which single ratio governs
the qualitative regime. Along the capital nullcline \eqref{eq:cnullcline},
equilibrium vitality is proportional to $\mu/\alpha$; the classification of
\S\ref{sec:equilibrium-lowvitality} therefore turns on $\mu/\alpha$ relative
to an explicit threshold. At baseline $\mu/\alpha=2.0$, against a flip value
of $2.163$ at which $V^{*}$ crosses $\Vstrat$---an $8\%$ margin. Across a
$\pm2\times$ log-uniform hypercube around the baseline, the probability that
a stable interior equilibrium classifies as low-vitality is approximately
$0.49$ (companion-code robustness sweep). The baseline is thus a genuinely
knife-edge illustration, and we present it as such.

\subsection{The low-vitality equilibrium and scope absorption}
\label{sec:equilibrium-lowvitality}

\begin{definition}[Strategic vitality threshold]
\label{def:threshold}
The strategic vitality threshold is $\Vstrat=\theta\,\Vmax$ with the
stipulated value $\theta=0.5$: the minimum vitality for sustained strategic
(as opposed to reactive) leadership work. The value of $\theta$ is a
labeling convention, not an estimated quantity; sensitivity of the
classification to $\theta$ is reported in \S\ref{sec:limitations}.
\end{definition}

\begin{definition}[Low-vitality equilibrium; outcome regimes]
\label{def:lowvitality}
A stable interior equilibrium $(V^{*},C^{*})$ is a \emph{low-vitality
equilibrium} if $V^{*}<\Vstrat$. A parameter vector is classified
\emph{sustainable} if its stable interior equilibrium has
$V^{*}\ge\Vstrat$, \emph{low-vitality} if $V^{*}<\Vstrat$, and
\emph{collapse-prone} if no stable interior equilibrium exists (equivalently
$\Cmax=0$ in the sense of Proposition~\ref{prop:carrying}). Earlier drafts
called the low-vitality equilibrium the ``zombie-leader equilibrium''; the
label is retired as a construct and survives only in legacy code names.
\end{definition}

At any interior equilibrium, \eqref{eq:cnullcline} gives
\begin{equation}
\label{eq:vstar}
V^{*} \;=\; \frac{\mu}{\alpha}\bigl(1+\varphi O^{*}\bigr),
\end{equation}
and at the baseline calibration the unique interior equilibrium is
\[
V^{*}=4.7025,\qquad C^{*}=32.0337,\qquad O^{*}=9.0084
\qquad(\varepsilon=0.1),
\]
a low-vitality equilibrium sitting $5.9\%$ below $\Vstrat$; the
$\varepsilon\to0$ closed form gives $V^{*}=4.67994$, $C^{*}=31.7325$
(Appendix~\ref{app:equilibrium}). Convergence is a damped spiral with
eigenvalues $-0.205\pm0.331i$ (Appendix~\ref{app:jacobian}).

\begin{corollary}[Scope absorption; intervention asymmetry]
\label{cor:scope}
In the model family \eqref{eq:dVdt}--\eqref{eq:complexity}, where the drain
kernel is the separable power law $\delta f(O)g(V)$ with $f(O)=O^{\gamma}$
and every appearance of $C$ in the flow is channeled through
$O=O_0+\beta C^{\eta}$, the equilibrium pair $(V^{*},O^{*})$ solves the
$\beta$-free system
\[
R\Bigl(1-\frac{V}{\Vmax}\Bigr)=\delta O^{\gamma}\frac{V}{V+\varepsilon},
\qquad
V=\frac{\mu}{\alpha}(1+\varphi O).
\]
Consequently, on the interior branch: (i) $V^{*}$ and $O^{*}$ are invariant
in $\beta$; (ii) the product $\beta\,C^{*\eta}=O^{*}-O_0$ is conserved. At
baseline ($\eta=1$), $\beta C^{*}=8.008$ at $\varepsilon=0.1$ and exactly
$7.933$ in the $\varepsilon\to0$ closed form---the invariant itself is
$\varepsilon$-dependent and must be quoted with its regularization
(Appendix~\ref{app:epsilon}).
\end{corollary}

\begin{proof}
Both equilibrium conditions are equations in $(V,O)$ alone; $\beta$ enters
the dynamics only through the map $C\mapsto O$. Given the unique root
$(V^{*},O^{*})$ (Theorem~\ref{thm:uniqueness}), inverting
\eqref{eq:complexity} yields $C^{*}=((O^{*}-O_0)/\beta)^{1/\eta}$, whence
$\beta C^{*\eta}=O^{*}-O_0$ independent of $\beta$.
\end{proof}

\begin{remark}[Honest status of Corollary~\ref{cor:scope}]
\label{rem:scope}
Earlier drafts headlined this result as a surprising discovery
(``Lemma~2''); the companion code retains that name. It is more honestly a
\emph{reparameterization invariance of the power-law kernel family}: with
$\eta=1$, $\beta$ multiplies a pure power of $C$ and is absorbed by
rescaling the units of capital, $C\mapsto\beta C$. It is \emph{not} a robust
law of leadership dynamics: under load or drain kernels in which the coupling
parameter also sets a saturation or curvature scale that cannot be absorbed
by rescaling $C$ (exponential and Hill specifications are the canonical
cases), the cancellation fails and $\partial V^{*}/\partial\beta\neq0$
generically (\S\ref{sec:robustness}, Appendix~\ref{app:bifurcation}). The
managerial payload survives within the declared family: \emph{scope reduction
alone cannot raise equilibrium vitality; recovery-side interventions---%
raising $R$, lowering $\mu$, raising $\alpha$---can.}
\end{remark}

\subsection{Stability: local, spectral, and global}
\label{sec:stability}

\begin{lemma}[Trace negativity: no Hopf bifurcation]
\label{lem:trace}
At any interior equilibrium of \eqref{eq:dVdt}--\eqref{eq:complexity}, for
all positive parameters, the Jacobian $J$ satisfies
$\operatorname{tr}J<0$ and $\det J>0$. Hence every interior equilibrium is
locally asymptotically stable, and no Hopf bifurcation occurs anywhere in
parameter space.
\end{lemma}

\begin{proof}[Proof sketch]
$J_{VV}=-R/\Vmax-\delta O^{\gamma}\varepsilon/(V+\varepsilon)^{2}<0$
unconditionally. For $J_{CC}$: the condition $\dot C=0$ at an interior
equilibrium forces $\alpha V/(1+\varphi O)=\mu$, so the first and third
terms of
$J_{CC}=\alpha V/(1+\varphi O)
-\alpha CV\varphi\,O'(C)/(1+\varphi O)^{2}-\mu$
cancel, leaving
$J_{CC}=-\alpha C V\varphi\,O'(C)/(1+\varphi O)^{2}<0$.
Thus $\operatorname{tr}J=J_{VV}+J_{CC}<0$. Moreover
$J_{VC}<0$ and $J_{CV}=\alpha C/(1+\varphi O)>0$, so
$\det J=J_{VV}J_{CC}-J_{VC}J_{CV}>0$. A Hopf bifurcation requires a purely
imaginary conjugate pair, i.e.\ $\operatorname{tr}J=0$, which is impossible;
this is consistent with (and implied more globally by) the Dulac certificate
$B=1/C$ of Theorem~\ref{thm:gas}. Full entries and baseline numerics are in
Appendix~\ref{app:jacobian}.
\end{proof}

\begin{proposition}[Carrying capacity as a flux threshold, not a bifurcation]
\label{prop:carrying}
Define the carrying capacity
\begin{equation}
\label{eq:cmax}
\Cmax \;=\;
\Biggl(\frac{(R/\delta)^{1/\gamma}-O_0}{\beta}\Biggr)^{1/\eta},
\end{equation}
the capital level at which the depletion ratio \eqref{eq:gamma-ratio}
reaches $\Gamma=1$ (evaluated at $V=\Vmax$); $\Cmax=44.99$ at baseline. Then
\begin{equation}
\label{eq:cmax-statics}
\frac{\partial \Cmax}{\partial\beta}<0,\qquad
\frac{\partial \Cmax}{\partial\delta}<0,\qquad
\frac{\partial \Cmax}{\partial R}>0 .
\end{equation}
$\Cmax$ is a \emph{flux threshold}: for $C>\Cmax$ the drain ceiling exceeds
the recovery ceiling and vitality is driven toward the depletion band
regardless of its current level. It is \emph{not} a bifurcation point: by
Lemma~\ref{lem:trace}, $\det J>0$ at every interior equilibrium for all
positive parameters, so no saddle-node (fold) bifurcation exists anywhere,
and no equilibria are created or destroyed at $\Cmax$. Earlier drafts'
description of $\Cmax$ as a ``critical threshold beyond which equilibria
cease to exist'' is retracted (Appendix~\ref{app:carrying}).
\end{proposition}

\begin{theorem}[Global asymptotic stability]
\label{thm:gas}
Let the interior equilibrium exist (Theorem~\ref{thm:uniqueness}) and define
the trapping rectangle $\Omega=[0,\Vmax]\times[0,\Ctrap]$ with
\begin{equation}
\label{eq:ctrap}
\Ctrap^{\eta} \;=\;
\frac{\bigl(\alpha\Vmax/\mu-1\bigr)/\varphi \;-\; O_0}{\beta},
\end{equation}
the unique capital level at which $V_c(O(C))=\Vmax$; $\Ctrap=102.67$ at
baseline. Then $\Omega$ is forward invariant, every trajectory with
$V_0\in[0,\Vmax]$, $C_0>0$ enters $\Omega$, and the interior equilibrium is
globally asymptotically stable on the \emph{open} set
$\{C>0\}\cap\Omega$. The set $\{C=0\}$ is invariant, and the axis
equilibrium $(V\approx9.934,\,C=0)$ is a saddle whose stable manifold lies
in $\{C=0\}$, so it attracts no interior trajectory.
\end{theorem}

\begin{proof}[Proof sketch]
(1)~\emph{Trapping:} on $C=\Ctrap$, $\alpha V/(1+\varphi O)\le
\alpha\Vmax/(1+\varphi O(\Ctrap))=\mu$, so $\dot C\le0$; the remaining edges
are handled by \S\ref{sec:invariance}. (2)~\emph{No closed orbits:} the Dulac
function $B(V,C)=1/C$ gives
$\nabla\!\cdot\!(B\mathbf{F})<0$ on $\{V>0,C>0\}$, term by term, ruling out
closed orbits and orbit graphics in the open quadrant (Bendixson--Dulac).
(3)~\emph{Limit sets:} $\Omega$ is compact; by Poincar\'e--Bendixson the
$\omega$-limit set of any interior trajectory is an equilibrium, and by
Theorem~\ref{thm:uniqueness}, Lemma~\ref{lem:trace}, and the saddle structure
of the axis equilibrium, that equilibrium is the interior one. The complete
argument, including the invariance verification edge by edge, is
Appendix~\ref{app:gas}.
\end{proof}

\begin{remark}[$\Ctrap\neq\Cmax$: the corrected trapping bound]
\label{rem:ctrap}
Earlier drafts used the rectangle $[0,\Vmax]\times[0,\Cmax]$ with
$\Cmax=44.99$. That rectangle \emph{leaks}: at $(V,C)=(\Vmax,\,44.99)$ one
computes $\dot C\approx+6.86>0$, and the trajectory from $(V_0,C_0)=(10,40)$
overshoots to $C\approx46.4$ before returning. The trapping bound
\eqref{eq:ctrap} and the carrying capacity \eqref{eq:cmax} are distinct
objects---a forward-invariance bound versus an energy-flux threshold---and
conflating them invalidated the earlier proof
(Appendix~\ref{app:gas}).
\end{remark}

\begin{remark}[Why Poincar\'e--Bendixson]
No elementary global Lyapunov function appears to exist for this system:
weighted quadratic forms centered at the equilibrium fail because the
cross terms generated by the $O^{\gamma}$ drain and the $1/(1+\varphi O)$
gate cannot be dominated globally. The Dulac-plus-Poincar\'e--Bendixson
route is therefore not a convenience but the appropriate tool for this
planar system.
\end{remark}

\subsection{Robustness of the specification}
\label{sec:robustness}

\paragraph{Linear drain ($\gamma=1$).}
With a linear complexity kernel the low-vitality regime disappears at
baseline: the unique stable equilibrium is \emph{sustainable}, at
$V^{*}\approx8.56$ ($C^{*}\approx83.5$). Nonlinear complexity scaling
($\gamma>1$) is therefore a necessary condition for the low-vitality
\emph{regime}---but not for the existence of an attractor, which persists at
$\gamma=1$. The qualitative theory stands or falls with the convexity of the
load--drain relationship, an empirically checkable premise.

\paragraph{Kernel dependence of scope absorption.}
Corollary~\ref{cor:scope} is a property of the power-law family. Under
exponential or Hill (saturating) kernels the $\beta$-invariance of $V^{*}$
fails, and $|\partial V^{*}/\partial\beta|$ grows with the departure from
separability (Proposition~P4 of \S\ref{sec:propositions};
Appendix~\ref{app:bifurcation}). We flag this boundary explicitly so the
intervention asymmetry is not over-generalized.
